% ---------------------------------------------------------------------------
% 論文中文草稿（NTCIR-19 AI CUP VeriPromiseESG Track 投稿用）
% 以 XeLaTeX 編譯：  xelatex paper_draft_zh_2026-08-30.tex
%
% 這份草稿的用途是決定「留什麼、怎麼說故事」。兩份進度報告合計約 3,100 行，
% 投稿只有 7 頁，所以本文的第一要務是取捨；附錄 A 記錄了每一個實驗的去留與理由。
% ---------------------------------------------------------------------------
\documentclass[11pt,a4paper]{ctexart}

% Overleaf：Menu -> Compiler 必須設為 XeLaTeX。
% 字體必須明確指定為繁體字集（TC）。若交給 ctex 自動挑選，它在 Overleaf 上會選到
% 簡體字型，繁體字會變成缺字方框。
\usepackage{xeCJK}
\setCJKmainfont{Noto Serif CJK TC}
\setCJKsansfont{Noto Sans CJK TC}
\setCJKmonofont{Noto Sans Mono CJK TC}

\usepackage{geometry}
\geometry{margin=2.3cm}
\usepackage{booktabs}
\usepackage{longtable}
\usepackage{array}
\usepackage{amsmath}
\usepackage[table]{xcolor}
\usepackage{enumitem}
\usepackage{hyperref}
\hypersetup{colorlinks=true, linkcolor=black, urlcolor=blue, citecolor=black}
\usepackage{tcolorbox}
\usepackage{tikz}
\usetikzlibrary{positioning}
\usepackage{fancyhdr}

\pagestyle{fancy}
\fancyhf{}
\fancyhead[L]{論文中文草稿}
\fancyhead[R]{\thepage}
\renewcommand{\headrulewidth}{0.4pt}

\newcommand{\fn}[1]{\texttt{#1}}

\title{\textbf{當判準沒有隨資料一起釋出}\\[0.3em]
       \large 一個 ESG 承諾驗證任務的診斷、可行的替代路徑，與棄權機制}
\author{以 VeriPromiseESG4K 為材料的獨立研究}
\date{2026 年 8 月 30 日\quad（草稿，用於決定取捨）}

\begin{document}
\ctexset{contentsname={目錄}, figurename={圖}, tablename={表},
         abstractname={摘要}, appendixname={附錄}}
\maketitle

\begin{abstract}
\noindent
判斷一段 ESG 永續報告書的佐證「夠不夠清楚」，是 VeriPromiseESG4K 這份繁體中文資料集
上最難的一個欄位。人類標註者在此欄的一致度高達 $\alpha = 0.940$，但已發表的語言模型
baseline 只有 macro-F1 0.3599，本文使用的微調集成系統也只有 0.4642。這個落差本身就是
問題：\textbf{人做得到、機器做不到，而且差距不隨模型變大而縮小}。本文因此不再嘗試
提高分數，而是問\textbf{困難的來源是什麼}。

\medskip
我們以四道相互獨立的檢驗排除了標註雜訊、模型能力與任務形式化三種解釋，並指出
真正的成因：\textbf{決定標籤的判準沒有隨資料一起釋出}。官方文件寫下的兩項判準
（措辭是否模糊、與承諾的連結是否明確）重現標註結果的能力只有 0.592 與 0.675 的
平衡準確率，低於官方\textbf{從未提及}的表面具體性線索（0.690）；官方唯一給出的
操作範例在 770 個四模型一致的段落中出現 0 次。要求四個前沿模型說明判斷依據後，
在\textbf{四個模型一致、卻與人類標註不同}的 149 個段落上，它們有 71.4\% 提到
「缺少數字」——一項沒有任何官方文件寫過的判準；而它們說某段「模糊」時，
84\% 的情況同一句話裡也在說缺少數字。這個代換在開源權重上複製成功，且在被逐條追問
四項官方判準時依然存在。

\medskip
在此診斷之上我們提出三項可行的做法：一個以\textbf{子句}為單位的跨度標註模型，把
一個從未被任何系統使用的標註資源變成可用任務（承諾抽取字元 F1 0.4434 $\to$ 0.5903）；
一個\textbf{棄權層}，使用不看模型內部的文字特徵，在 80\% 覆蓋率下把整列聯合正確率
從 0.4880 提升到 0.5325；以及一個\textbf{異質組合}，把偏誤強度不同的模型平均，使
準確率從 0.8052 提升到 0.8300。我們並以\textbf{成員之間的錯誤相關係數}解釋為什麼
同樣是集成，四個前沿模型投票無效（$\phi = +0.498$）而上述異質組合有效
（$\phi = +0.307$）。

\medskip
\textbf{關鍵詞}：ESG、承諾驗證、標註判準、選擇性預測、集成學習、繁體中文
\end{abstract}

\tableofcontents
\newpage

% ===========================================================================
\section{緒論}

上市公司每年出版永續報告書，其中充滿「我們將於 2030 年達成碳中和」這類承諾。有些
附有具體佐證，有些只是空泛宣示——後者正是漂綠（greenwashing）。VeriPromiseESG4K
把這個問題形式化為四個段落層級的分類欄位：這段有沒有承諾、有沒有佐證、\textbf{佐證
的品質如何}、以及該承諾何時可被驗證。

\medskip
四個欄位的難度極不平均。人類標註者在四個欄位上的一致度都在 0.94 以上，但自動方法在
\textbf{佐證品質}這一個欄位上遠遠落後，且這個落差在已發表的 baseline、我們的系統、
以及本文測試的四個前沿商用模型上一致出現。本文處理的問題不是「怎麼再高一點」，而是：

\begin{tcolorbox}[colback=blue!5, colframe=blue!55]
\textbf{這個子任務的困難來自哪裡？它與其他三個子任務的困難是同一種嗎？如果不是，
還有什麼是做得到的？}
\end{tcolorbox}

\paragraph{本文的貢獻}
\begin{enumerate}[leftmargin=2.2em, itemsep=3pt]
  \item \textbf{一個可檢驗的診斷}。以四道獨立檢驗排除標註雜訊、模型能力、任務形式化
        三種解釋，並以模型自述的判斷依據直接證實機制。診斷對其他三個子任務做出預測，
        該預測被驗證（第 \ref{sec:diag} 節）。
  \item \textbf{一個可行的替代任務}。資料集提供覆蓋率 100\% 的跨度標註，從未被任何
        已發表的系統使用。我們指出正確的文字單元是\textbf{子句}而非句子，
        並建成模型（第 \ref{sec:tagger} 節）。
  \item \textbf{兩個把診斷變成系統行為的機制}：棄權與異質組合（第 \ref{sec:act} 節）。
  \item \textbf{一個對集成何時有效的解釋}：以成員錯誤相關係數說明為什麼前沿模型
        投票無效而異質組合有效（第 \ref{sec:voting} 節）。
\end{enumerate}

% ===========================================================================
\section{任務與資料}

VeriPromiseESG4K 含 4,000 個段落，取自 49 家台灣上市公司的 2024 年永續報告書，
切分為訓練 1,000／驗證 1,000／測試 2,000。每段有四個分類欄位與\textbf{兩個跨度欄位}。

\begin{table}[h]
\centering\small
\begin{tabular}{@{}clp{5.6cm}rr@{}}
\toprule
 & 欄位 & 問題 & 標註者一致度 $\alpha$ & 我們的 macro-F1 \\
\midrule
\textcircled{1} & \fn{promise\_status} & 這段有沒有 ESG 承諾？ & 0.977 & 0.8191 \\
\textcircled{2} & \fn{evidence\_status} & 承諾有沒有具體佐證？ & 0.962 & 0.7132 \\
\rowcolor{yellow!15}
\textcircled{3} & \fn{evidence\_quality} & 佐證品質如何？ & 0.940 & \textbf{0.4642} \\
\textcircled{4} & \fn{verification\_timeline} & 何時可被驗證？ & 0.955 & 0.6545 \\
\midrule
 & \fn{promise\_string} & 承諾是哪一句（跨度） & --- & 未評分 \\
 & \fn{evidence\_string} & 佐證是哪一句（跨度） & --- & 未評分 \\
\bottomrule
\end{tabular}
\caption{四個分類欄位與兩個未被使用的跨度欄位。$\alpha$ 為 Krippendorff's Alpha，
資料集發布者報告的標註者一致度。最右欄是本文系統在\textbf{全部 2,000 段、四個類別
（含 \fn{N/A}）}上的成績，用於呈現四個欄位之間的落差；後文所有比較改用另一把尺，
見第 \ref{sec:scale} 節。}
\end{table}

\subsection{本文使用的系統，以及三個一再出現的名詞}
\label{sec:system}

本文多處以「本文的系統」為比較對象。為了讓後續每一張表都讀得懂，先把三個名詞說清楚。

\paragraph{編碼器（encoder）} 一個把中文文字轉成數值向量的預訓練神經網路。它事先在
大量中文語料上訓練過，因此已經「認得中文」，但不知道本任務。要它做本任務，需在其
輸出端接一個小型分類層，再用本資料集的標註資料繼續訓練——這個步驟稱為\textbf{微調}
（fine-tuning）。本文使用的三個編碼器各約\textbf{3 億個參數}，屬於中型模型；作為
對照，本文比較的商用前沿模型參數量通常大兩到三個數量級。

\paragraph{五折交叉驗證（5-fold cross-validation）} 把 2,000 列訓練資料\textbf{平均
切成五份}，每次拿其中四份訓練、剩下一份預測，輪流五次。好處有二：每一列都得到一個
「由沒看過它的模型」所做的預測（本文稱之為\textbf{折外預測}，用於誠實地挑選門檻與
權重）；以及最後會得到\textbf{五個}訓練資料略有不同的模型。

\paragraph{軟投票（soft voting）} 把多個模型的預測合起來。之所以稱「軟」，是因為
平均的是每個模型輸出的\textbf{類別機率}，而不是它們各自的最終答案（後者稱硬投票）。
平均機率保留了「模型有多猶豫」這個資訊。

\begin{table}[h]
\centering\small
\begin{tabular}{@{}ll@{}}
\toprule
成員 & 三種中文預訓練編碼器：\fn{chinese-roberta-wwm-ext-large}、\\
     & \fn{chinese-macbert-large}、\fn{Erlangshen-DeBERTa-v2-320M} \\
訓練 & 每個編碼器各做一次五折交叉驗證，因此共有 $3 \times 5 = 15$ 個微調後的模型 \\
推論 & 15 個模型對每個測試段落各輸出一組類別機率，取平均後選機率最大的類別 \\
\bottomrule
\end{tabular}
\caption{本文使用的系統。後文所有「本文的系統」「我們微調的編碼器」皆指此。}
\end{table}

\paragraph{什麼叫「同一協定」} 本文多處說某兩個實驗「使用同一編碼器、同一五折協定」，
意思是：\textbf{兩者的編碼器、資料切分方式、隨機種子、訓練輪數、批次大小與學習率
全部相同，只有被比較的那一項不同}。這是為了確保觀察到的差異來自那一項，而不是來自
訓練設定的差異。

\paragraph{為什麼報 macro-F1} 本文沿用 macro-F1，因為它是類別不平衡下的標準做法，
也是這份資料集既有 baseline 使用的指標，兩者因此可以直接對照。所謂 macro-F1，是
\textbf{先算每個類別各自的 F1，再取未加權平均}——不論該類別有幾筆資料。

\medskip
這個定義在本資料集上有一個必須先講明的後果：\textcircled{3} 的 \fn{Misleading} 類別
在全部 4,000 段中\textbf{只有 5 筆}（訓練 1、驗證 1、測試 3），卻與有 1,106 筆的
\fn{Clear} 同等計入平均。因此本文多處會看到\textbf{同一個改動讓使用者看到的結果變好、
macro-F1 卻下降}；每次出現時我們都會同時報告準確率，讓兩者分開判讀。

\subsection{全文只用一把尺}
\label{sec:scale}

上表最右欄是在\textbf{全部 2,000 段、四個類別}上算的，因為那是四個欄位唯一的共同
基礎。但本文後續所有比較都聚焦在\textbf{佐證品質}這一個欄位，而該欄位真正要判斷的
只有\textbf{有佐證的那些段落}——沒有佐證的段落其答案必為 \fn{N/A}，是由前一個欄位
決定的，混進來只會稀釋要研究的東西。

\medskip
因此\textbf{第 \ref{sec:diag} 節之後的每一個數字，都在同一把尺上}：

\begin{table}[h]
\centering\small
\begin{tabular}{@{}ll@{}}
\toprule
評測列 & 測試集中 \fn{evidence\_status}=\fn{Yes} 且品質標籤有效的 \textbf{1{,}335 段} \\
類別 & \fn{Clear}（1{,}106）、\fn{Not Clear}（226）、\fn{Misleading}（3），共三類 \\
指標 & macro-F1 為主，準確率並列 \\
\bottomrule
\end{tabular}
\end{table}

\paragraph{本文用到的評測名詞} 以下五個名詞在後文反覆出現，一次講完。

\begin{table}[h]
\centering\small
\begin{tabular}{@{}p{2.6cm}p{12cm}@{}}
\toprule
準確率 & 猜對的段落佔全部的比例。直觀，但在類別不平衡時會被多數類別灌水。 \\
\addlinespace
macro-F1 & 先算每個類別各自的 F1，再取\textbf{未加權}平均，不論該類別有幾筆。 \\
\addlinespace
平衡準確率 & 先算每個類別各自的正確率，再取平均。與 macro-F1 精神相同，但更容易解讀：
\textbf{隨機亂猜的分數等於「1 除以類別數」}。三個類別時是 0.333，兩個類別時是 0.500。 \\
\addlinespace
AUC & 用來衡量一個分數「排序」的能力：把段落依該分數由高到低排，真正該排前面的有多常
排在前面。0.5 等於亂排，1.0 是完美。 \\
\addlinespace
95\% 信賴區間 & 以 bootstrap（把測試集重複隨機重抽數千次，每次重算一遍）估計出來的
不確定範圍。\textbf{若區間跨越 0，就無法排除「其實沒有差別」}。 \\
\bottomrule
\end{tabular}
\end{table}

\paragraph{換尺之後有一個結果會反過來，必須說明} 本文的三編碼器五折集成在這把尺上是
macro-F1 \textbf{0.4339}（準確率 0.7618），\textbf{低於}單一編碼器五折的 0.4655
（準確率 0.8052）。原因不是集成比較差，而是\textbf{兩者被問的問題不同}：集成必須自己
判斷這一段有沒有佐證，因此在這 1,335 段中有 115 段輸出了 \fn{N/A}（一律計為錯）；
單一編碼器則是在「已知這段有佐證」的前提下，只從三個類別中選一個。

\medskip
後文以\textbf{單一編碼器五折}作為本文系統的代表，因為它與被比較的語言模型問的是同一
個問題。集成的 0.4642 只用於上表，不與後文任何數字並排。

% ===========================================================================
\section{診斷：困難的來源}
\label{sec:diag}

\subsection{論證結構}

我們以四道檢驗逐一排除替代解釋。圖 \ref{fig:arg} 是整條論證的結構。

\begin{figure}[h]
\centering
\begin{tikzpicture}[
  % anchor=north：每格由上緣往下長，行數不同也不會往上撞到鄰列
  box/.style={draw, rounded corners=2pt, align=center, font=\small,
              inner sep=4pt, text width=32mm, anchor=north},
  ar/.style={->, gray}]

\node[box, fill=blue!10, draw=blue!50, text width=64mm] (q) at (5.85,6.2)
  {\textbf{為什麼子任務\textcircled{3}這麼難？}};

\node[box, fill=gray!12]  (a) at (0,    5.0) {解釋一\\標籤有雜訊};
\node[box, fill=gray!12]  (b) at (3.9,  5.0) {解釋二\\模型能力不足};
\node[box, fill=gray!12]  (c) at (7.8,  5.0) {解釋三\\任務形式化錯誤};
\node[box, fill=green!14, draw=green!50!black] (d) at (11.7,5.0)
  {\textbf{解釋四}\\\textbf{判準沒有隨資料釋出}};

\node[box, fill=gray!4]   (a2) at (0,   3.0)
  {\textbf{排除}\\標註者一致度\\$\alpha = 0.940$};
\node[box, fill=gray!4]   (b2) at (3.9, 3.0)
  {\textbf{排除}\\四個前沿模型\\0.354--0.449\\低於本文的 0.466};
\node[box, fill=gray!4]   (c2) at (7.8, 3.0)
  {\textbf{排除}\\六種輸入差異\\皆在 $\pm0.006$ 內\\而對照組 $-0.0185$};
\node[box, fill=green!8, draw=green!50!black] (d2) at (11.7,3.0)
  {\textbf{成立}\\官方判準重現正解\\僅 0.592／0.675\\低於未提及的 0.690};

\draw[ar] (q.south) -- (a.north); \draw[ar] (q.south) -- (b.north);
\draw[ar] (q.south) -- (c.north); \draw[ar] (q.south) -- (d.north);
\draw[ar] (a.south) -- (a2.north); \draw[ar] (b.south) -- (b2.north);
\draw[ar] (c.south) -- (c2.north); \draw[ar] (d.south) -- (d2.north);
\end{tikzpicture}
\caption{論證結構。三種替代解釋各以一項獨立量測排除，第四種被四項證據支持。}
\label{fig:arg}
\end{figure}

\subsection{檢驗一與二：不是雜訊，也不是能力}

\paragraph{不是標籤有雜訊} 資料集發布者報告的標註者一致度為 $\alpha = 0.940$。也就是說
受過訓練的人類在這個欄位上高度一致——如果標籤只是雜訊，人類彼此也不會一致。困難因此
不能歸因於「這題本來就沒有答案」。

\paragraph{不是模型能力不足} 我們以三家廠商的四個前沿商用模型在同一批 1,335 段上
作答，不提供任何範例。

\begin{table}[h]
\centering\small
\begin{tabular}{@{}lrrr@{}}
\toprule
系統 & macro-F1 & 準確率 & 輸出 \fn{N/A} 的次數 \\
\midrule
\fn{claude-opus-5} & 0.3534 & 0.5708 & 21 \\
\fn{claude-sonnet-5} & 0.4487 & 0.7835 & 11 \\
\fn{gpt-5.6-sol} & 0.4339 & 0.7551 & 45 \\
\fn{gemini-3.7-flash} & 0.4178 & 0.7543 & 88 \\
\fn{Qwen3-8B}（開源，零樣本） & 0.3914 & 0.7700 & --- \\
\midrule
\rowcolor{green!12}
\textbf{本文微調的編碼器}（3 億參數） & \textbf{0.4655} & \textbf{0.8052} & --- \\
\bottomrule
\end{tabular}
\end{table}

\textbf{參數量大兩到三個數量級的模型，沒有一個比得上那個 3 億參數、但用本資料集微調過
的編碼器。} 困難因此不是「模型不夠強」。

\paragraph{這個比較有一處不對稱，必須寫明} 四個前沿模型的提示詞給的是\textbf{四個}
選項（含 \fn{N/A}），本文的編碼器只從\textbf{三個}類別中選。前者因此可能把答案花在
一個在這批列上必定錯誤的選項上——實際發生了 11 到 88 次，佔 0.8\% 至 6.6\%，一律計為
錯誤。這對前沿模型不利，但幅度不足以翻轉結論：最強的
\fn{claude-sonnet-5} 只誤用該選項 11 次，其 0.4487 與本文的 0.4655 之間的差距與此
無關。

\subsection{檢驗三：不是任務形式化}

第三種可能是\textbf{題目問錯了}：也許不該把整段丟給模型，而該先找出「佐證是哪一句」，
再只對那一句判斷品質。資料集正好提供了這個跨度標註，所以這個重構做得到。

\medskip
我們把\textbf{餵給分類器的文字}換成六種，其餘一切維持不變（同一個編碼器、同一組五折
切分、同一組超參數，見第 \ref{sec:system} 節）：

\begin{table}[h]
\centering\small
\begin{tabular}{@{}clrl@{}}
\toprule
條件 & 輸入 & macro-F1 & 對 A 的差（95\% 信賴區間）\\
\midrule
A & 完整段落（現行做法） & 0.4655 & --- \\
B & 標註提供的佐證跨度 & 0.4600 & $-0.0056\ [-0.0226, +0.0114]$ \\
C & 標註提供的承諾＋佐證 & 0.4695 & $+0.0039\ [-0.0120, +0.0193]$ \\
\rowcolor{yellow!15}
D & 等長隨機片段（對照） & 0.4470 & $\mathbf{-0.0185\ [-0.0389, -0.0003]}$ \\
E & \textbf{預測的}佐證跨度 & 0.4665 & $+0.0009\ [-0.0103, +0.0124]$ \\
F & \textbf{預測的}承諾＋佐證 & 0.4704 & $+0.0048\ [-0.0069, +0.0168]$ \\
\bottomrule
\end{tabular}
\caption{重構任務沒有幫助，而這次的實驗有足夠檢定力證明。}
\end{table}

關鍵不在跨度條件的數字，而在\textbf{對照組 D 的信賴區間不跨零}——實驗偵測得到 0.018
量級的差異；而所有跨度條件與 A 的差都在 $\pm 0.006$ 以內，信賴區間全部跨零。這不是
「量不出來」，是「效果不存在」。

\paragraph{一個免費的副產物} 條件 F 使用模型\textbf{自己預測}的跨度，成績不低於讀
整段。一個先找出承諾句與佐證句再判斷品質的系統，分類準確度沒有損失，卻多輸出兩個
對稽核者有用的欄位。

\subsection{一批特別有用的段落：四個模型一致、卻與人不同}
\label{sec:split}

檢驗四與後續分析都用到同一批段落，其來歷如下。四個前沿模型對同一批段落各自作答之後，
我們把它們\textbf{答案完全相同}的段落挑出來，再依那個共同答案與人類標註是否一致，
分成兩組。

\begin{table}[h]
\centering\small
\begin{tabular}{@{}clr@{}}
\toprule
步驟 & 條件 & 剩下幾段 \\
\midrule
1 & 測試集全部 & 2{,}000 \\
2 & 其中有佐證、且品質標籤有效（第 \ref{sec:scale} 節那把尺） & 1{,}335 \\
3 & 四個模型\textbf{答案完全相同} & 770 \\
4a & \quad ……而且與人類標註相同 $\rightarrow$ \textbf{對照組} & 621 \\
\rowcolor{yellow!15}
4b & \quad ……但與人類標註\textbf{不同} $\rightarrow$ \textbf{爭議組} & \textbf{149} \\
\bottomrule
\end{tabular}
\caption{兩組段落的來歷。四個模型彼此不一致的 565 段不列入——那是模型自己就沒共識
的段落，無法用來談「模型與人的分歧」。}
\end{table}

\paragraph{為什麼這兩組特別有用} 兩組都是四個獨立系統一致同意的段落，差別只在
\textbf{人同不同意}。因此把兩組拿來對比時，「模型有多確定」這個因素已經被固定住了，
剩下的差異就是\textbf{人機分歧本身}。

\paragraph{分歧是單向的} 爭議組的 149 段中，人標 \fn{Clear} 的有 119 段（80\%），而
四個模型一致說 \fn{Not Clear} 的有 120 段。也就是說模型不是隨機出錯，而是系統性地
\textbf{比人更嚴格}。這個方向性是後續所有分析的起點。

\subsection{檢驗四：判準沒有隨資料釋出}

官方對 \fn{evidence\_quality} 的說明只有兩處，且都在資料欄位說明中：要求評估
「語意明確」「易於理解」「未使用模糊字眼（例如：持續推進、努力改善）」，以及
「與承諾的連結是否明確、細節是否充足」。\textbf{兩項判準都沒有提到數字。}

我們把可能的判準分成三組，各自寫成可以機械計算的特徵，再看哪一組最能重現什麼。

\paragraph{三組特徵，以及它們各自的出處} 必須先講明一件事：\textbf{三組特徵全都是
本研究寫出來的}。三者的差別不在誰寫的，而在\textbf{它們試圖操作化的那條判準，官方
有沒有寫下來}。

\begin{table}[h]
\centering\small
\begin{tabular}{@{}p{0.6cm}p{7.6cm}p{6.2cm}@{}}
\toprule
組 & 特徵內容（全部由本研究撰寫） & 它操作化的判準 \\
\midrule
A & 官方逐字給的兩個模糊詞，加上本研究另行擴充的 22 個同語域詞彙，及其密度 &
    「未使用模糊字眼」\newline \textbf{官方寫下的} \\
\addlinespace
B & 承諾句與佐證句的字元重疊率、佐證中不在承諾裡的比例、佐證長度、長度比 &
    「與承諾的連結是否明確、細節是否足夠」\newline \textbf{官方寫下的} \\
\addlinespace
\rowcolor{yellow!12}
C & 年份、數字、百分比、計量單位、金額、涉及方指認、第三方查證，共 7 條正規表示式，
    加上數字密度 & 「有沒有具體數字」\newline \textbf{官方從未提過的} \\
\bottomrule
\end{tabular}
\end{table}

\paragraph{怎麼比} 每一組特徵各自餵進一個邏輯迴歸，做五折交叉驗證，用折外預測算
\textbf{平衡準確率}。目標有兩個：一次是重現\textbf{人類標註}，一次是重現\textbf{四個
模型的共識答案}。目標都化約為二元（\fn{Clear} 對 \fn{Not Clear}，因 \fn{Misleading}
僅 3 筆不足以擬合），故隨機水準為 0.500，評測列為第 \ref{sec:scale} 節那 1,335 段中
標籤為此二類的 1,332 段。

\begin{table}[h]
\centering\small
\begin{tabular}{@{}llrr@{}}
\toprule
組 & 它操作化的判準 & 重現人類標註 & 重現模型共識 \\
\midrule
A & 措辭模糊（官方寫下的） & 0.592 & 0.654 \\
B & 與承諾的連結（官方寫下的） & 0.675 & 0.708 \\
\rowcolor{yellow!15}
C & 有沒有具體數字（\textbf{官方從未提過}） & \textbf{0.690} & \textbf{0.786} \\
\bottomrule
\end{tabular}
\caption{平衡準確率，五折折外，隨機水準 0.500。\textbf{官方寫下的兩條判準，是最不能
重現官方標註結果的兩組。}}
\end{table}

再加兩項事實：官方唯一給出的操作範例（「持續推進」「努力改善」）在 770 個段落中
\textbf{出現 0 次}；三組判準\textbf{全都}比較能解釋模型而非標註者。

\subsection{機制的直接證據：問模型自己}

以上都是從數字反推模型在看什麼。更直接的做法是\textbf{問模型自己}：我們要求四個
模型在給標籤的同時，用一句話說明「你是看了段落中的什麼才這樣判斷」，再把這些理由
自動歸類。

\paragraph{歸類可不可信} 歸類是用正規表示式（文字比對規則）做的，所以必須先驗證它。
我們隨機抽出 60 筆理由人工逐筆判讀，與程式的結果比對，一致程度以 \textbf{Cohen's
Kappa} 衡量——這個統計量把「碰巧一致」扣掉之後才計分，0 代表等同亂猜、1 代表完全
一致，一般認為 0.8 以上屬高度一致。承載本文論證的兩個類別分別是 \textbf{0.963}
（措辭模糊）與 \textbf{0.861}（缺少數字）；其餘類別偏低，因此只列出供參考，不用於
任何主張。

\paragraph{表中的數字是什麼} 每一格是\textbf{該組段落中，模型的理由提到這件事的
比例}（對照組 200 段、爭議組 149 段，每個模型各一次；\fn{gemini-3.7-flash} 有 9 筆
逾時，已排除）。最右欄的 $p$ 值經 \textbf{Holm 校正}——同時檢定多個項目時，光憑運氣
就會有幾項看起來顯著，Holm 校正把這個機率壓回去，因此校正後仍顯著的項目才可信。

\medskip
$^{\dagger}$\textbf{這一列不能用來下結論。} 人工驗證顯示，這個類別的自動歸類
Cohen's Kappa 為 \textbf{0.000}——在 60 筆抽樣中人工判定它出現 0 次、程式誤判 3 次。
也就是說這個偵測器\textbf{偵測不到東西}，因此「不顯著」與「偵測不到」在此無法區分。
本文\textbf{不從這一列推論}模型有沒有使用官方的連結判準；該判準的地位由前一節的
表（B 組 0.675）認定。

\begin{table}[h]
\centering\small
\begin{tabular}{@{}lrrr@{}}
\toprule
模型自述的判斷依據 & 對照組 & 爭議組 & Holm 校正 $p$ \\
\midrule
官方判準：措辭模糊 & 0.112 & 0.719 & $1.7\times10^{-117}$ \\
\rowcolor{gray!12}
官方判準：與承諾的連結$^{\dagger}$ & 0.077 & 0.094 & 0.25（不顯著）\\
\rowcolor{yellow!15}
\textbf{官方從未提及：缺少具體數字} & 0.149 & \textbf{0.714} & $2.9\times10^{-100}$ \\
引用「有具體數據」 & \textbf{0.439} & 0.088 & $1.2\times10^{-45}$ \\
\bottomrule
\end{tabular}
\end{table}

「措辭模糊」與「缺少數字」的比率幾乎相同，這不是巧合：\textbf{模型說某段「模糊」時，
同一句話裡有 84.1\% 同時在說「缺少數字」}。模型在\textbf{形式上}遵守了官方寫下的
判準，卻把「模糊」\textbf{實例化成「沒有數字」}。

\paragraph{兩項排除} 這個代換\textbf{不是提示詞格式造成的}：把官方四項判準逐字拆成
四題、強迫模型分別作答，那個共現率只從 0.876 降到 0.808（此處為單一模型
\fn{claude-sonnet-5} 的數值，因為只有它在三種問法下的作答分布保持穩定，可跨問法比較）；
而單純要求模型「一步一步想再回答」（\textbf{chain-of-thought}，逐步推理）反而讓共現率
\textbf{上升}到 0.905。給更多推理空間，買到的是更有把握地套用同一條規則。

\medskip
\textbf{它也不是某一家實驗室的產物。} 我們把完全相同的提示詞與完全相同的 349 段
（149 爭議 + 200 對照）交給一個\textbf{開源的 80 億參數模型}（\fn{Qwen3-8B}）——規模
小一個量級、由第四家機構訓練、權重公開可查，而且它\textbf{沒有參與}第 \ref{sec:split}
節那兩組的挑選。結果同樣的不對稱重現了：

\begin{table}[h]
\centering\small
\begin{tabular}{@{}lrrr@{}}
\toprule
\fn{Qwen3-8B} 自述的判斷依據 & 對照組 & 爭議組 & 差 \\
\midrule
缺少具體數字 & 0.055 & 0.315 & $+0.260$ \\
引用「有具體數據」 & 0.300 & 0.081 & $-0.219$ \\
\bottomrule
\end{tabular}
\caption{欄位數值為\textbf{該組段落中，模型的理由提到這件事的比例}。例如 0.315 表示
在 149 段爭議段落裡，有 31.5\% 的理由提到「缺少數字」。}
\end{table}

在它說某段「模糊」的理由裡，有 71.9\% 同一句話也在說缺少數字（四個前沿模型是
84.1\%）。

\subsection{診斷對其他子任務的預測，以及驗證}

若診斷正確，它應該能說明\textbf{哪些子任務會容易}。官方文件對另外三個子任務都給了
可在文字表面檢查的線索（\textcircled{1} 五個表示承諾的詞、\textcircled{2} 數據／方法／
措施、\textcircled{4} 以 2024 年為起算點的明確區間），只有 \textcircled{3} 給的是形容詞。

\begin{table}[h]
\centering\small
\begin{tabular}{@{}lrrrr@{}}
\toprule
 & \textcircled{1} & \textcircled{2} & \textcircled{3} & \textcircled{4} \\
\midrule
官方字面文字能重現正解（隨機校正後） & $+0.298$ & $+0.375$ & $\mathbf{+0.002}$ & $+0.210$ \\
補上一個可用操作化所需的增量 & $+0.132$ & $-0.003$ & $\mathbf{+0.375}$ & $+0.048$ \\
\midrule
表面特徵解釋「系統」勝過解釋「標註者」 & $+0.041$ & $+0.171$ & $\mathbf{+0.246}$ & $\mathbf{+0.001}$ \\
\bottomrule
\end{tabular}
\end{table}

\begin{tcolorbox}[colback=green!5, colframe=green!45, title=兩個軸，而非一個]
\textcircled{3} 的官方字面文字重現正解只有 $+0.002$——\textbf{等同隨機}。

\medskip
但 \textcircled{4} 才是關鍵對照組：它一樣難，\textbf{分歧卻是零}——系統與標註者讀著
同一批表面證據，一起失敗。那是\textbf{普通的困難}，而普通的困難對普通的補救有反應：
資料集發布者自己的檢索增強 baseline 把 \textcircled{4} 從 0.3090 拉到 0.4227，把 \textcircled{3}
只從 0.3549 拉到 0.3599。

\medskip
\textbf{\textcircled{3} 不是最難的子任務，而是唯一一個困難來自「對問題的分歧」而非
「對答案的分歧」的子任務。}
\end{tcolorbox}

% ===========================================================================
\section{可行的替代路徑：子句層級的跨度標註}
\label{sec:tagger}

資料集為每個正例標出承諾句與佐證句，覆蓋率 100\%。但既有的評測只針對四個分類欄位，
因此\textbf{據我們所知沒有任何已發表的系統使用過這兩欄}。

\subsection{關鍵在選對文字單元}

\begin{table}[h]
\centering\small
\begin{tabular}{@{}lrr@{}}
\toprule
把跨度化約為單元的邊界對齊率 & 句子為單元 & \textbf{子句為單元} \\
\midrule
標註邊界與單元邊界吻合 & 24\% & \textbf{71\%} \\
化約後可達的字元 F1 上限 & 0.5414 & \textbf{0.9748} \\
\bottomrule
\end{tabular}
\end{table}

中文子句以逗號分隔，而標註者的切點大量落在逗號上。改以子句為單元後，任務從困難的
「自由文字位置抽取」變成標準的「逐子句二元判斷」。

\subsection{怎麼算對錯}

跨度抽取的答案是「一段文字」而不是一個標籤，因此需要另一組指標。設模型畫出來的字元
集合為 $P$、人類標註畫出來的為 $G$：

\begin{table}[h]
\centering\small
\begin{tabular}{@{}p{2.6cm}p{12cm}@{}}
\toprule
字元 F1 & 把 $P$ 與 $G$ 當成兩袋字元，計算重疊程度。畫太多會扣分（精確率低），畫
太少也會扣分（召回率低）。 \\
\addlinespace
完全匹配 & $P$ 與 $G$ \textbf{一字不差}的段落比例。這是最嚴格的指標，也是實務上
真正有用的——差一個字，抽出來的句子就不是那一句。 \\
\bottomrule
\end{tabular}
\end{table}

兩者要一起看。把\textbf{整段照抄}當答案，字元 F1 可以拿到 0.6004（因為佐證通常佔段落
一半以上），但完全匹配只有 0.4\%——那是一個字元 F1 看起來不差、實際上毫無用處的策略。

\subsection{模型與結果}

整段編碼一次，再依斷詞器提供的字元位置把每個子句的向量讀出來，逐子句判斷「這個子句
是不是承諾（或佐證）的一部分」。

\begin{table}[h]
\centering\small
\begin{tabular}{@{}lrrrr@{}}
\toprule
系統 & 承諾：字元 F1 & 完全匹配 & 佐證：字元 F1 & 完全匹配 \\
\midrule
整段照抄 & 0.4183 & 5.5\% & 0.6004 & 0.4\% \\
含線索詞的句子（最強啟發式） & 0.4434 & 11.9\% & 0.4004 & 1.7\% \\
Qwen3-8B 零樣本 & 0.2704 & 1.5\% & 0.5537 & 1.4\% \\
Qwen3-8B 四個範例 & 0.3038 & 9.5\% & 0.4727 & 6.8\% \\
\rowcolor{green!12}
\textbf{子句標註模型} & \textbf{0.5903} & \textbf{14.0\%} & \textbf{0.6666} & \textbf{11.8\%} \\
\bottomrule
\end{tabular}
\caption{語言模型即使拿到官方逐字定義與四個範例，仍低於「整段照抄」；4-shot 之下
還有 518 個回傳片段是段落裡不存在的文字。}
\end{table}

\subsection{消融：哪一個設計在起作用}

三個種子的標準差為 0.0033（承諾）與 0.0010（佐證），故雜訊帶為 $\pm 0.0066$／$\pm 0.0020$。

\begin{table}[h]
\centering\small
\begin{tabular}{@{}lrr@{}}
\toprule
拿掉一個設計選擇 & 承諾：對完整模型 & 佐證：對完整模型 \\
\midrule
\rowcolor{yellow!15}
改用句子單元 & $\mathbf{-0.1032}$ & $\mathbf{-0.0307}$ \\
\rowcolor{yellow!15}
移除段落脈絡（遮蔽跨子句注意力） & $\mathbf{-0.0590}$ & $\mathbf{-0.0238}$ \\
換成 macbert 編碼器 & $-0.0066$ & $+0.0013$ \\
\bottomrule
\end{tabular}
\caption{兩個形式化選擇遠超雜訊帶，編碼器選擇落在帶內。成績來自任務的形式化，
不是來自挑到好的預訓練模型。}
\end{table}

執行成本：單張 RTX 4090，五折訓練 296 秒，每千個子句推論 0.31 秒。

\begin{tcolorbox}[colback=blue!5, colframe=blue!50]
\textbf{為什麼這個任務成功，而子任務\textcircled{3}失敗？} 同一個理論解釋兩者：
跨度抽取的判準\textbf{就在文字表面}（措辭特徵），子任務\textcircled{3}的判準不在。
\end{tcolorbox}

% ===========================================================================
\section{把診斷變成系統行為}
\label{sec:act}

\subsection{棄權：知道何時不該回答}

若一部分段落無法由系統拿得到的文字決定，正確的行為是\textbf{棄權}——不猜，交給人。
這只有在\textbf{不看正解就能認出那些段落}時才有用，而第 \ref{sec:diag} 節的判準特徵
正是從段落本身算出來的，沒有用到答案。

\paragraph{怎麼評量一個會棄權的系統} 用\textbf{覆蓋率}（coverage）：系統選擇回答的
段落佔全部的比例。覆蓋率 100\% 表示全部都答，80\% 表示挑最有把握的八成回答、其餘
交給人。棄權一定會讓剩下的準確率上升——因為丟掉的是難題——所以有意義的比較不是
「棄權有沒有用」，而是\textbf{在同一個覆蓋率下，不同的棄權訊號誰挑得準}，以及誰
勝過\textbf{隨機丟棄}這個對照。

\begin{table}[h]
\centering\small
\begin{tabular}{@{}lrrrr@{}}
\toprule
棄權訊號 & 預測「會錯」AUC & 覆蓋 90\% & 覆蓋 80\% & 覆蓋 50\% \\
\midrule
模型自己的信心（標準做法） & 0.7688 & 0.8303 & 0.8614 & 0.9251 \\
\rowcolor{yellow!12}
\textbf{判準特徵}（不看模型內部） & \textbf{0.7937} & 0.8369 & 0.8680 & 0.9476 \\
\rowcolor{green!12}
\textbf{兩者合併} & \textbf{0.8191} & \textbf{0.8453} & \textbf{0.8820} & \textbf{0.9626} \\
隨機（對照） & 0.4499 & 0.7978 & 0.7949 & 0.7799 \\
\bottomrule
\end{tabular}
\caption{子任務\textcircled{3}的準確率對覆蓋率。不看模型內部的文字特徵，比模型
自己的信心更會預測它要犯的錯，而且兩者互補。}
\end{table}

推廣到\textbf{整列}（四欄同時正確）時，另有一個零成本訊號：預測若違反標註階層
（沒有承諾則其餘三欄必為 \fn{N/A}），該列\textbf{必定}不可能四欄全對——因為正解一律
符合階層。測試集有 89 列（4.5\%）如此。

\begin{table}[h]
\centering\small
\begin{tabular}{@{}lrrrr@{}}
\toprule
覆蓋率 & 判準特徵 & 結構不一致 & 兩者合併 & 隨機 \\
\midrule
100\% & 0.4880 & 0.4880 & 0.4880 & 0.4880 \\
80\% & 0.5100 & 0.5056 & \textbf{0.5325} & 0.4881 \\
50\% & 0.5660 & 0.5110 & \textbf{0.5830} & 0.4920 \\
\bottomrule
\end{tabular}
\caption{整列聯合正確率對覆蓋率。}
\end{table}

\paragraph{另一種處理方式，以及它的適用範圍} 對這些列還有一個直覺的做法：把它們
\textbf{投影}到合法的標籤組合上（例如既然說沒有承諾，就把其餘三欄改成 \fn{N/A}）。
在本文的多任務架構上，這麼做會扣 0.0130 的四欄平均 macro-F1——它把預測推向 \fn{N/A}，
以稀有類別的召回率換取多數類別的精確度。

\medskip
\textbf{但這個代價不是普遍的。} 換成四個欄位各自獨立訓練的專用模型架構，同樣的投影
代價是 $+0.0013$ 與 $+0.0003$——等於零，方向還相反。機制說得通：多任務架構共用編碼器，
四個分類頭高度相關，違反邏輯通常代表整列的表徵都錯了，投影只是把錯誤傳播出去；四個
獨立模型各自失準時，違反邏輯較常是「一個頭錯、三個對」，那正是投影修得好的情況。

\begin{tcolorbox}[colback=gray!8, colframe=gray!55]
\textbf{因此投影的代價是「多任務架構」的性質，不是指標或任務的性質}；本文不把它列為
可推廣的結論。而\textbf{棄權}不需要在架構之間做這個區分：它移除的是\textbf{本來就
不可能正確}的列，這一點對任何架構都成立。
\end{tcolorbox}

\subsection{異質組合：平均進一個偏誤較弱的模型}

\begin{table}[h]
\centering\small
\begin{tabular}{@{}lrrl@{}}
\toprule
組合 & macro-F1 & 準確率 & 準確率變化（bootstrap 95\% 信賴區間）\\
\midrule
微調 transformer 單獨 & 0.4655 & 0.8052 & --- \\
Qwen3-8B 零樣本單獨 & 0.3914 & 0.7700 & $-0.0352\ [-0.0614, -0.0105]$ \\
\rowcolor{green!12}
折外擬合權重 0.45 & \textbf{0.4742} & \textbf{0.8300} & $\mathbf{+0.0247\ [+0.0090, +0.0390]}$ \\
\bottomrule
\end{tabular}
\caption{折外搜尋把權重放在 0.45，而且是放在一個\textbf{每一欄都比 transformer 差}
的模型上。這只有在兩者錯的列不同時才划算。}
\end{table}

% ===========================================================================
\section{討論：集成何時有效}
\label{sec:voting}

本文出現了兩個方向相反的集成結果，值得放在一起解釋。四個前沿模型投票\textbf{沒有
增益}；一個微調 transformer 與一個零樣本開源模型的組合\textbf{有增益}。集成理論對此
有現成的答案——平均只在成員\textbf{錯得不一樣}時有價值——而本文剛好可以直接量那個
量：兩個系統的錯誤向量之間的相關係數 $\phi$。圖~\ref{fig:vote} 是本文量到的
兩個點，以及兩個尚未量測的位置。

\begin{figure}[h]
\centering
\begin{tikzpicture}[
  b/.style={draw, rounded corners=2pt, align=center, font=\small,
            inner sep=3pt, text width=31mm},
  ar/.style={->, gray}]
\node[b, fill=gray!6, draw=gray!45] (f5) at (0,0)
  {五折之間\\\textit{$\phi$ 未量測}};
\node[b, fill=gray!6, draw=gray!45] (bb) at (3.8,0)
  {三種編碼器之間\\\textit{$\phi$ 未量測}\\增益 $+0.0142$};
\node[b, fill=red!8] (fr) at (7.6,0)
  {\textbf{四個前沿模型}\\$\phi = +0.498$\\\textbf{無增益}};
\node[b, fill=green!14] (het) at (11.4,0)
  {\textbf{異質組合}\\$\phi = +0.307$\\準確率 $+0.0247$};
\draw[ar] (f5) -- (bb); \draw[ar] (bb) -- (fr); \draw[ar] (fr) -- (het);
\node[font=\small\itshape] at (5.7,-2.2)
  {成員錯誤相關係數 $\phi$ 愈低，平均愈有價值};
\end{tikzpicture}
\caption{本文量到的兩個點（右二），以及尚未量測的兩個（左二，灰）。灰色兩格的
$\phi$ 需要保留每一折、每個 backbone 的個別預測才能計算，本文未做。}
\label{fig:vote}
\end{figure}

\begin{table}[h]
\centering\small
\begin{tabular}{@{}lrr@{}}
\toprule
成員組合 & 錯誤相關 $\phi$ & 平均之後 \\
\midrule
\rowcolor{red!8}
四個前沿模型彼此（平均） & $+0.498$ & 無增益（先前實驗七個模型全負） \\
\rowcolor{green!12}
我們的 transformer 與 Qwen3-8B & $\mathbf{+0.307}$ & \textbf{準確率} $\mathbf{+0.0247}$ \\
\fn{claude-opus-5} 與 Qwen3-8B（最低的一對） & $+0.051$ & 未測 \\
\bottomrule
\end{tabular}
\caption{錯誤相關係數解釋了哪些集成有用。四個前沿模型「彼此的一致度（0.765）高於
它們與標註者的一致度（0.701）」——它們共享同一個表面具體性代換，所以一起錯。}
\end{table}

\begin{tcolorbox}[colback=gray!8, colframe=gray!55, title=能主張的與不能主張的]
\textbf{能主張的}：本文的兩個集成結果——前沿模型投票無增益、異質組合有增益——由成員
錯誤相關係數解釋，而該係數的差異（$+0.498$ 對 $+0.307$）又由本文的診斷解釋：四個
前沿模型共享同一個表面具體性代換，所以一起錯。這對「該找什麼樣的成員」給出一個具體
的建議——\textbf{不是更大的模型，而是偏誤強度不同的模型}。

\medskip
\textbf{不能主張的}：本文\textbf{沒有}量測五折之間或三種 backbone 之間的錯誤相關，
因此不能宣稱已經把本文系統內部的五折軟投票與上述結果統一成一個連續的解釋。「平均在成員錯得
不一樣時才有價值」也是集成學習的既有結論，不是本文的發現；本文的貢獻只在於
\textbf{指出這個資料集上共享偏誤的具體內容}，使該係數變得可以事先預測。
\end{tcolorbox}

% ===========================================================================
\section{結論與限制}

子任務 \textcircled{3} 的困難不在標籤雜訊、模型能力或任務形式化，而在\textbf{決定
標籤的判準沒有隨資料一起釋出}。資料集論文自述，標註準則是在每週檢討會議中反覆修訂、
並由永續領域專家裁定後才達到 $\alpha=0.940$——那個過程的產物從未公開。文字傳不了
這條判準；\textbf{標註範例可以}——用本資料集的 1,345 筆標註資料去微調一個開源 8B
模型，它在這個欄位上的準確率從 0.770 升到 0.840（附錄 B）。同一條判準，寫成文字傳不
過去，做成範例就傳得過去。

\paragraph{限制}
\begin{itemize}[leftmargin=2em, itemsep=2pt]
  \item \fn{Misleading} 全資料集僅 5 筆，任何關於該類別的結論都受此限制。
  \item 模型理由的自動編碼由本文作者撰寫並由同一作者人工驗證；應由未參與的標註者
        重做。
  \item 錯誤相關與增益的關係目前只有兩個可用的點，尚不足以稱為曲線；五折之間與
        三種 backbone 之間的錯誤相關未量測，需保留每個成員的個別預測才能補上。
  \item 所有結果限於繁體中文 ESG 報告書與此單一資料集。
\end{itemize}

% ===========================================================================
\appendix
\section{每個實驗的去留}

兩份進度報告合計約 3,100 行、46 支程式。下表是投稿取捨。

\begin{longtable}{@{}clp{6.6cm}c@{}}
\toprule
程式 & 內容 & 在論文中的角色 & 去留 \\
\midrule
\endhead
26 & 前沿模型天花板 & 檢驗二（排除能力） & \textbf{留} \\
29 & 分歧結構 0.765 vs 0.701 & 檢驗四動機、第 \ref{sec:voting} 節 & \textbf{留} \\
33 & 分歧特徵顯著性（Holm） & 檢驗四（支撐） & 留（併表）\\
34 & 三組判準比較 & \textbf{檢驗四核心} & \textbf{留} \\
35 & 模型自述判斷依據 & \textbf{機制直證} & \textbf{留} \\
39 & 跨子任務對照 & 診斷的預測與驗證 & \textbf{留} \\
40 & 判準拆解、CoT 對照 & 排除「提示詞格式」解釋 & 留（一段）\\
43 & 開源模型複製 & 排除「單一實驗室」解釋 & 留（一句）\\
28 & 子句 vs 句子上限 & 第 \ref{sec:tagger} 節動機 & \textbf{留} \\
32 & 子句標註模型 & \textbf{正面結果一} & \textbf{留} \\
42 & 消融、三種子、runtime & 場次形式要求 & \textbf{留} \\
37 & 語言模型跨度 baseline & 第 \ref{sec:tagger} 節對照 & \textbf{留} \\
36 & 六種輸入條件 & 檢驗三（排除形式化） & \textbf{留} \\
38 & 選擇性預測 & \textbf{正面結果二} & \textbf{留} \\
46 & 整列棄權 & 正面結果二（推廣） & \textbf{留} \\
45 & 異質組合、公平比較 & \textbf{正面結果三}、第 \ref{sec:voting} 節 & \textbf{留} \\
05,14 & 強制一致的代價 & 棄權那節的對照；架構性質，不列為通則 & 留（一段）\\
\midrule
44 & Self-REF 信心符號 & 假設被推翻；支持「範例可傳達」 & 附錄 \\
41 & 編碼人工驗證 & 限制段 & 附錄 \\
\midrule
07 & TF-IDF 跨度探針 & 已被 36 取代 & 不留 \\
09 & 語意關聯度（假設 C） & 已被 34/35 涵蓋 & 不留 \\
11 & Qwen 零樣本 & 已被 26/43 涵蓋 & 不留 \\
19 & LLM 資料擴增 & 負面且無獨立價值 & 不留 \\
21,22 & LLM 投票權重 & 併入第 \ref{sec:voting} 節一句 & 不留 \\
23--25 & 檢索增強 & 引用資料集既有 baseline 即可 & 不留 \\
27 & 跨度清理重做 & 已被 36 涵蓋 & 不留 \\
30,31 & 人工裁決、鄰近一致性 & 訊號弱（7 點差距） & 不留 \\
\bottomrule
\end{longtable}

\paragraph{取捨原則} 留下的每一項都在論證鏈上有\textbf{不可替代}的位置：排除一個
替代解釋、支撐一個正面結果、或滿足形式上的完整性。九項失敗介入中只有「強制一致」
留下，因為只有它產生了可推廣的方法學結果；其餘八項在正文中以一句話帶過。

\section{文字傳不了的判準，範例傳得過去}
\label{app:selfref}

正文第 \ref{sec:diag} 節的結論是：官方文件的\textbf{文字}不足以把這個欄位的判準傳達
給任何模型。一個自然的追問是——那什麼傳得過去？

\medskip
我們用本資料集訓練＋驗證集中有佐證的 \textbf{1{,}345 筆標註資料}，去微調一個開源的
80 億參數模型（\fn{Qwen3-8B}）。這裡的「微調」不是要它學會回答，而是把\textbf{標註者
的判斷本身}當成範例餵給它。

\begin{table}[h]
\centering\small
\begin{tabular}{@{}lrr@{}}
\toprule
 & 微調前 & 微調後 \\
\midrule
在 1{,}335 段測試資料上的準確率 & 0.770 & \textbf{0.840} \\
\bottomrule
\end{tabular}
\end{table}

\begin{tcolorbox}[colback=blue!5, colframe=blue!50]
同一條判準：\textbf{寫成文字傳不過去}（把定義寫進提示詞、把四項判準逐條追問，都失敗
——正文第 \ref{sec:diag} 節）；\textbf{做成範例就傳得過去}（1,345 筆標註，準確率
$+0.070$）。

\medskip
這也解釋了正文一開始那個乍看奇怪的現象：一個 3 億參數的微調編碼器，為什麼會在同一批
段落上勝過四個參數量大兩三個數量級的前沿模型。差別不在模型大小，而在
\textbf{有沒有看過這批標註者的判斷}。
\end{tcolorbox}

\paragraph{附帶的一個否定結果} 我們同時依 Self-REF（Chuang 等人，ICML 2025）的做法，
在模型的答案後面加上兩個特殊符號 \fn{<CN>}（有把握）與 \fn{<UN>}（沒把握），訓練它在
自己答對時吐前者、答錯時吐後者，藉此讀出一個信心分數。本研究原本預測\textbf{這種
信心訊號抓不到判準不匹配的錯誤}——因為在爭議段落上模型並非不會，而是有把握地套用
另一條判準。\textbf{這個預測是錯的}：它在爭議段落上仍達到 0.754 的 AUC。原因就在上表
——微調已經把一部分判準傳給了模型，模型的不確定性因此跟著那條判準走。

\end{document}
